\documentclass[hyphens,aspectratio=169,dvipsnames]{beamer}
\usepackage{graphicx}
\usepackage{xcolor}
\usepackage{amssymb}
\usepackage{amsmath}
\usepackage{mathtools}
\usepackage{textcomp}
\usepackage{moresize}
\usepackage{framed}
\usepackage{minted}
\usepackage{relsize}
\usepackage{tikz}
\usepackage{comment}
\usetikzlibrary{shapes.geometric, arrows, positioning}
\usetheme{Berlin}
\usecolortheme[RGB={0,95,47}]{structure}
\beamertemplatenavigationsymbolsempty

\makeatletter
\AtBeginEnvironment{minted}{\dontdofcolorbox}
\def\dontdofcolorbox{\renewcommand\fcolorbox[4][]{##4}}
\makeatother

\newcommand{\textpf}[1]{\texttt{\color{black}\fcolorbox{lightgray}{lightgray}{#1}}}

\begin{document}

\begin{frame}

    \textbf{Multi-file C programs}

    \begin{itemize}
        \item C files/modules: *.c
        \item Header files: *.h
    \end{itemize}
    
\end{frame}

\begin{frame}[fragile]{Multi-File C Programs}
    \begin{itemize}
        \pause \item Recall \textpf{my\_log2} and \textpf{my\_abs} from last class.
        \pause \item Suppose that we want to use them in the following program:
    \end{itemize}
    \inputminted{c}{main.c}

    % verbally explain what the ul suffix on an integer literal means.
    \pause \begin{center} \textbf{Task}: Copy down this program as-is into a file named \textpf{main.c} and attempt to compile it. \end{center}
\end{frame}


\begin{frame}[fragile]{Understanding the Errors}
    \scriptsize
    \begin{verbatim}
main.c: In function ‘main’:
main.c:2:12: error: implicit declaration of function ‘my_log2’ [-Wimplicit-function-declaration]
    2 |     return my_log2(my_abs(-100));
      |            ^~~~~~~
main.c:2:20: error: implicit declaration of function ‘my_abs’ [-Wimplicit-function-declaration]
    2 |     return my_log2(my_abs(-100));
      |                    ^~~~~~
    \end{verbatim}
\end{frame}

\begin{frame}[fragile]{What is Implicit Function Declaration?}
    \pause \begin{center} The compiler considers a function \textpf{f} to be ``implicitly declared" when \textpf{f} is called before the compiler knows that \textpf{f} exists. \end{center}
\end{frame}

\begin{frame}[fragile]{Fixing Implicit Function Declaration}
    \begin{center} How do we fix this? \end{center}
    \begin{itemize}
        \pause \item Copy and paste the definitions of \textpf{my\_log2} and \textpf{my\_abs} into \textpf{main.c}.
            \begin{itemize}
                \pause \item Doesn't scale well. We can't just cram everything into a single file!
            \end{itemize}
        \pause \item Define the functions in another file, then instruct the compiler to read both files.
    \end{itemize}
\end{frame}

\begin{frame}[fragile]{Activity \#1.0: Building a Multi-File Program}
    \begin{enumerate}
        \pause \item Take the definitions of \textpf{my\_log2} and \textpf{my\_abs}, and place them in a new file named \textpf{my\_math.c}.
        \pause \item Attempt to compile both \textpf{my\_math.c} and \textpf{main.c} together by passing them both to the compiler driver.
        \begin{itemize}
            \pause \item Why doesn't this work?
            \pause \item Hint: The compiler considers a function  \textpf{f} to ``implicitly declared" when \textpf{f} is called before the compiler knows that \textpf{f} exists.
        \end{itemize}
    \end{enumerate}
\end{frame}

\begin{frame}[fragile]{Activity \#1.1: Telling the Compiler to Shut Up}
    \begin{itemize}
        \pause \item The compiler still fails with the same error, because under the hood, files are compiled one at a time, and then linked together.
        \pause \item Thus, when the compiler is building \textpf{main.c}, it still doesn't know about the functions in \textpf{my\_math.c}.
        \pause \item You can tell the compiler to stop worrying about this by passing the \textpf{-Wno-implicit-function-declaration} flag on the command line.
        \pause \item This works because the linker resolves function references, and the linker has access to the results of compiling both files.
    \end{itemize}
\end{frame}

\begin{frame}[fragile]{Activity \#1.1: Telling the Compiler to Shut Up}
    \begin{itemize}
        \item Remember from last week that the linker links together multiple (compiled) object files into an executable file.
        \item In order to do so, it needs to resolve function references (this also applies to some globally defined variables).
    \end{itemize}

        Example: \\
        \{file1.c; file2.c\} -> compilation -> \{file1.o; file2.o\} -> linking -> \{my\_prog\}
\end{frame}



\begin{frame}[fragile]{Activity \#1.2: Why \textpf{-Wno-implicit-function-declaration} is Bad}
    \begin{center} \textbf{NEVER USE \textpf{-Wno-implicit-function-declaration}!!!} \end{center}

    \begin{itemize}
        \pause \item The compiler doesn't know the argument counts and types for implicitly-declared functions.
        \pause \item The linker doesn't care about argument counts and types at all.
        \pause \item Try modifying \textpf{main.c} to call the functions with the wrong numbers of arguments. What happens?
    \end{itemize}
\end{frame}

\begin{frame}[fragile]{Activity \#1.3: Making the Compiler Happy}
    \begin{itemize}
        \pause \item To inform the compiler about your math functions in \textpf{main.c}, those functions must be declared.
        \pause \item A function declaration is just like a function definition, but no code block is provided.
        \pause \item For example, the declaration for \textpf{my\_log2} might look like this:
            \begin{minted}{c}
unsigned char my_log2(unsigned long x);
            \end{minted}
        \pause \item Modify \textpf{main.c} to add declarations for your math functions before the definition of \textpf{main}. This should resolve the compiler errors without using the evil flag.
    \end{itemize}
\end{frame}

\begin{frame}[fragile]{Activity \#1.4: Your First Header File}
    \begin{itemize}
        \pause \item This is still kind of ugly. We don't want to litter our programs with a million function declarations.
        \pause \item To avoid this, function declarations are typically placed in their own files, called header files.
        \pause \item Make a header file, \textpf{my\_math.h}, and move the declarations from \textpf{main.c} into it.
        \pause \item Then, instruct the preprocessor to paste its contents into \textpf{main.c} by sticking \textpf{\#include "my\_math.h"} at the top of \textpf{main.c}.
    \end{itemize}
\end{frame}

\begin{frame}[fragile]{\textpf{\#include}}
    \begin{itemize}
        \pause \item Recall that the C preprocessor evaluates special directives within your C code.
        \pause \item \textpf{\#include} is one such directive.
        \pause \item The \textpf{\#include} directive tells the preprocessor to paste the contents of another file into the current file.
        \pause \item It takes a single argument, which is a path surrounded by either \textpf{""} or \textpf{<>}.
        \pause \item If the path is surrounded by \textpf{""}, the path is relative to the current directory. % Technically, it just tries relative to the current directory first, then falls back to the system include path. Too much work to explain in a slide though.
        \pause \item If the path is surrounded by \textpf{<>}, the path is relative to the system's default include path (\textpf{/usr/include}).
    \end{itemize}
\end{frame}

\begin{frame}[fragile]{\textpf{struct}s}
    \begin{itemize}
        \pause \item C provides a few keywords for defining new data types.
        \pause \item \textpf{struct} is one such keyword.
        \pause \item The \textpf{struct} keyword allows for the definition of ``record types," (aka ``product types").
        \pause \item Basically, a \textpf{struct} is a type that is defined as a bundle of some other types.
    \end{itemize}
\end{frame}

\begin{frame}[fragile]{\textpf{struct} Definition}
    Here's an example of a \textpf{struct} definition:
    \begin{minted}{c}
struct example_struct {
    unsigned char first_field;
    int second_field;
    short third_field;
};
    \end{minted}
\end{frame}

\begin{frame}[fragile]{Instantiating a \textpf{struct}}
    A \textpf{struct} can be instantiated in three ways:
    \begin{itemize}
        \pause \item Uninitialized
        \pause \item Using an initializer list
        \pause \item Using a designated initializer
    \end{itemize}
\end{frame}

\begin{frame}[fragile]{Uninitialized \textpf{struct} Instantiation}
    \begin{itemize}
        \pause \item To instantiate a \textpf{struct example\_struct} named \textpf{s} without initializing its members, do this:
            \begin{minted}{c}
struct example_struct s;
            \end{minted}
        \pause \item Note that accessing uninitialized data is undefined behavior!
    \end{itemize}
\end{frame}

\begin{frame}[fragile]{\textpf{struct} Initializer Lists}
    \begin{itemize}
        \pause \item To instantiate a \textpf{struct example\_struct} named \textpf{s} using an initializer list, do this:
            \begin{minted}{c}
struct example_struct s = {255, -1, 1000};
            \end{minted}
        \pause \item This sets \textpf{first\_field} to 255, \textpf{second\_field} to -1, and \textpf{third\_field} to 1000.
        \pause \item If you use an initializer list with too few fields, the remaining fields are default-initialized (set to 0).
    \end{itemize}
\end{frame}

\begin{frame}[fragile]{Designated \textpf{struct} Initialization}
    \begin{itemize}
        \pause \item To instantiate a \textpf{struct example\_struct} named \textpf{s} using a designated initializer, do this:
            \begin{minted}{c}
struct example_struct s = {
    .first_field = 255,
    .second_field = -1,
    .third_field = 1000
};
            \end{minted}
        \pause \item This sets \textpf{first\_field} to 255, \textpf{second\_field} to -1, and \textpf{third\_field} to 1000.
        \pause \item Any fields not populated in the designated initializer are default-initialized.
    \end{itemize}
\end{frame}

\begin{frame}[fragile]{\textpf{struct} Field Access}
    \begin{itemize}
        \pause \item To access in the fields in a \textpf{struct}, use the \textpf{.} operator.
        \pause \item For example, the following code makes an uninitialized \textpf{struct example\_struct}, then populates its fields:
            \begin{minted}{c}
struct example_struct s;
s.first_field = 255;
s.second_field = -1;
s.third_field = 1000;
            \end{minted}
    \end{itemize}
\end{frame}

\begin{frame}[fragile]{Activity \#2: Fractions!}
    \begin{itemize}
        \pause \item Write a header file, \textpf{fractions.h}, that defines a \textpf{struct} that represents a fraction, and declares functions for fraction multiplication, division, addition, and subtraction.
        \pause \item Write a C file, \textpf{fractions.c}, that defines the functions declared in \textpf{fractions.h}.
        \pause \item Don't worry about keeping fractions in simplest form.
    \end{itemize}
\end{frame}

\end{document}
